\subsection{Đánh giá mô hình nhận diện ký hiệu}

Mô-đun nhận diện là tầng đầu tiên của toàn bộ pipeline nên mọi sai số ở đây đều có thể lan xuống mô-đun
ghép câu và mô-đun dịch. Vì mô hình hiện tại là một mô hình pre-trained (TFLite từ cuộc thi Kaggle
ASL Signs), cách đánh giá sẽ khác so với một mô hình tự huấn luyện: nhóm không có quá trình
training/validation thông thường, mà tập trung vào đánh giá chất lượng suy luận thực tế.

\subsubsection{Các metric có thể đánh giá}

Mặc dù mô hình nhận diện trong đề tài là mô hình pre-trained, mô-đun này vẫn cần được đánh giá như
một bộ phân loại đa lớp. Điểm khác biệt là nhóm không báo cáo loss huấn luyện hay validation loss do
không huấn luyện lại mô hình; thay vào đó, phần đánh giá tập trung vào chất lượng dự đoán trên tập
kiểm thử gán nhãn và tốc độ suy luận. Các metric phù hợp bao gồm:
\begin{itemize}
    \item \textbf{Top-1 Accuracy}: tỉ lệ dự đoán đúng nhãn ký hiệu trên tập kiểm thử;
    \item \textbf{Top-3/Top-5 Accuracy}: tỉ lệ nhãn đúng nằm trong top-k dự đoán;
    \item \textbf{Macro Precision, Macro Recall, Macro F1}: đánh giá cân bằng theo từng lớp, tránh việc
    một vài lớp xuất hiện nhiều làm sai lệch kết quả chung;
    \item \textbf{Weighted F1}: F1 có trọng số theo số mẫu của từng lớp;
    \item \textbf{Confusion matrix}: ma trận nhầm lẫn giữa các lớp ký hiệu;
    \item \textbf{Average inference time}: thời gian suy luận trung bình trên mỗi chuỗi landmarks;
    \item \textbf{FPS}: số khung hình xử lý được trong một giây, quan trọng cho demo realtime;
    \item \textbf{Đánh giá định tính}: quan sát trực tiếp khi sử dụng webcam.
\end{itemize}

\begin{table}[H]
\centering
\caption{Ý nghĩa của các metric cho mô-đun nhận diện}
\label{tab:recognition-metric-meanings}
\begin{tabularx}{\textwidth}{|p{3cm}|X|X|}
\hline
\textbf{Metric} & \textbf{Ý nghĩa} & \textbf{Cách diễn giải} \\
\hline
Top-1 Accuracy & Tỉ lệ mẫu có lớp dự đoán cao nhất trùng đúng nhãn thật. & Càng cao càng tốt; là chỉ số chính của bài toán phân loại. \\
\hline
Top-3 Accuracy & Tỉ lệ mẫu mà nhãn thật xuất hiện trong 3 dự đoán cao nhất. & Hữu ích khi số lớp lớn và nhiều dấu hiệu gần nhau. \\
\hline
Top-5 Accuracy & Tỉ lệ mẫu mà nhãn thật xuất hiện trong 5 dự đoán cao nhất. & Cho thấy mô hình có ``nhìn ra'' vùng ứng viên đúng hay không. \\
\hline
Macro F1 & Trung bình F1 theo từng lớp ký hiệu. & Phù hợp khi muốn xem mô hình có công bằng giữa các lớp hay không. \\
\hline
Confusion matrix & Bảng đếm nhãn thật và nhãn dự đoán. & Dùng để phát hiện các cặp ký hiệu dễ bị nhầm. \\
\hline
Latency / FPS & Tốc độ suy luận trên mỗi mẫu hoặc theo khung hình. & Quyết định khả năng dùng trong realtime demo. \\
\hline
\end{tabularx}
\end{table}

Với tập kiểm thử gồm $N$ mẫu, nhãn thật $y_i$ và danh sách dự đoán top-k là $\hat{Y}_{i,k}$, các chỉ số
Top-k Accuracy được tính như sau:
\[
    TopK = \frac{1}{N}\sum_{i=1}^{N}\mathbf{1}(y_i \in \hat{Y}_{i,k})
\]

Precision, Recall và F1 theo từng lớp được tính từ số lượng true positive ($TP$), false positive ($FP$)
và false negative ($FN$):
\[
    Precision = \frac{TP}{TP + FP}, \quad
    Recall = \frac{TP}{TP + FN}, \quad
    F1 = \frac{2 \times Precision \times Recall}{Precision + Recall}
\]

\subsubsection{Bối cảnh mô hình pre-trained}

Cần nhấn mạnh rằng mô hình TFLite trong dự án là giải pháp đạt giải nhất cuộc thi Kaggle
\textit{Google -- Isolated Sign Language Recognition}. Cuộc thi này có hơn 1.000 đội tham gia
và đánh giá trên dữ liệu kiểm thử chính thức của Kaggle. Vì vậy, chất lượng của mô hình đã được
kiểm chứng ở quy mô cạnh tranh quốc tế.

Tuy nhiên, kết quả trong cuộc thi Kaggle được đo trên dữ liệu landmarks đã chuẩn bị sẵn, trong khi
pipeline thực tế của dự án phải trích xuất landmarks từ webcam theo thời gian thực. Điều này có thể
tạo ra chênh lệch giữa chất lượng offline (trên dữ liệu Kaggle) và chất lượng online (trên webcam
thực tế), do:
\begin{itemize}
    \item chất lượng camera, ánh sáng và nền ảnh hưởng đến độ chính xác của MediaPipe;
    \item tốc độ ký hiệu và khoảng cách tới camera có thể khác so với điều kiện trong dữ liệu
    huấn luyện;
    \item một số từ vựng trong 250 lớp có thể gần nhau về hình thức, dễ gây nhầm.
\end{itemize}

\begin{table}[H]
\centering
\caption{Thông tin tham chiếu từ cuộc thi Kaggle ASL Signs}
\label{tab:kaggle-asl-reference}
\begin{tabularx}{\textwidth}{|X|c|}
\hline
\textbf{Thuộc tính} & \textbf{Giá trị} \\
\hline
Tên cuộc thi & Google -- Isolated Sign Language Recognition \\
\hline
Số lớp từ vựng & 250 \\
\hline
Nguồn mô hình & Giải pháp đạt giải nhất (1st place solution) \\
\hline
Định dạng & TensorFlow Lite (.tflite) \\
\hline
Kích thước mô hình & $\sim$11 MB \\
\hline
\end{tabularx}
\end{table}

\subsubsection{Công cụ đánh giá định lượng}

Để phần nhận diện không chỉ dừng ở mức ``chạy mô hình pre-trained'', nhóm bổ sung script
\codepath{external/asl_baseline/evaluate_asl_250.py}. Script này có hai chế độ:
\begin{itemize}
    \item \textbf{benchmark-only}: đo thời gian suy luận của mô hình TFLite trên tensor landmarks giả lập
    có đúng kích thước đầu vào \texttt{N x 543 x 3};
    \item \textbf{manifest evaluation}: đọc file CSV gồm hai cột \texttt{path,label}, chạy suy luận trên
    các file landmarks \texttt{.npy}/\texttt{.npz}, sau đó xuất Top-1, Top-3, Top-5 Accuracy, Macro
    Precision, Macro Recall, Macro F1, Weighted F1, latency và confusion matrix.
\end{itemize}

Lệnh benchmark tốc độ đã dùng trong snapshot này:
\begin{verbatim}
venv\Scripts\python.exe external\google-asl-250\evaluate_asl_250.py ^
  --benchmark-only --runs 30 --warmup 5
\end{verbatim}

Khi có tập kiểm thử landmarks gán nhãn, lệnh đánh giá phân loại có dạng:
\begin{verbatim}
venv\Scripts\python.exe external\google-asl-250\evaluate_asl_250.py ^
  --manifest test_manifest.csv --output-dir external\google-asl-250\eval_results
\end{verbatim}

\begin{table}[H]
\centering
\caption{Kết quả benchmark suy luận TFLite của mô-đun nhận diện}
\label{tab:recognition-latency-benchmark}
\begin{tabularx}{\textwidth}{|c|c|c|c|}
\hline
\textbf{Độ dài chuỗi} & \textbf{Mean latency (ms)} & \textbf{P95 latency (ms)} & \textbf{FPS tương đương} \\
\hline
16 frame & 13,26 & 16,57 & 75,40 \\
\hline
32 frame & 26,26 & 36,91 & 38,09 \\
\hline
60 frame & 44,92 & 49,92 & 22,26 \\
\hline
\end{tabularx}
\end{table}

Các số liệu trong Bảng~\ref{tab:recognition-latency-benchmark} chỉ đo phần suy luận TFLite trên CPU,
chưa bao gồm thời gian đọc webcam, tiền xử lý ảnh và trích xuất landmarks bằng MediaPipe. Tuy vậy, kết
quả này cho thấy riêng mô hình phân loại đủ nhanh để phục vụ demo realtime, vì trong code hiện tại mô
hình chỉ được gọi mỗi 8 frame thay vì gọi ở mọi frame.

\begin{table}[H]
\centering
\caption{Trạng thái các metric phân loại của mô-đun nhận diện trong snapshot hiện tại}
\label{tab:recognition-classification-status}
\begin{tabularx}{\textwidth}{|p{4cm}|p{3.2cm}|X|}
\hline
\textbf{Metric} & \textbf{Trạng thái} & \textbf{Ghi chú} \\
\hline
Top-1 Accuracy & Chưa công bố số & Cần tập kiểm thử landmarks/video có nhãn thật. Repository hiện chỉ có mô hình pre-trained và label map, chưa có test manifest cho recognition. \\
\hline
Top-3 / Top-5 Accuracy & Chưa công bố số & Script đã hỗ trợ tính tự động khi có dữ liệu gán nhãn. \\
\hline
Macro Precision / Recall / F1 & Chưa công bố số & Không thể tính đúng nếu thiếu \texttt{y\_true}. Việc tự đưa số khi chưa có tập test sẽ không có giá trị học thuật. \\
\hline
Confusion matrix & Chưa công bố bảng & Sẽ được xuất ra file \texttt{recognition\_confusion\_matrix.csv} khi chạy đánh giá với manifest gán nhãn. \\
\hline
Latency / FPS & Đã đo & Có số liệu benchmark TFLite trong Bảng~\ref{tab:recognition-latency-benchmark}. \\
\hline
\end{tabularx}
\end{table}

\subsubsection{Đánh giá thực nghiệm qua demo realtime}

Vì snapshot hiện tại chưa có tập test nhận diện riêng, phần đánh giá thực nghiệm qua webcam được dùng
để kiểm tra khả năng vận hành thực tế của pipeline. Đây không thay thế cho Top-k Accuracy hay F1, nhưng
có giá trị bổ sung vì pipeline thật phải đi qua camera, MediaPipe và cửa sổ trượt thời gian. Các tiêu
chí đánh giá bao gồm:
\begin{itemize}
    \item mô hình có nhận đúng các từ vựng phổ biến như \texttt{hello}, \texttt{drink},
    \texttt{food}, \texttt{please}, \texttt{thankyou} hay không;
    \item mô hình phản hồi nhanh bao lâu sau khi người dùng bắt đầu ký hiệu;
    \item mô hình có ổn định (không nhảy liên tục giữa các nhãn) khi người dùng giữ tay
    trong tư thế;
    \item mô hình xử lý thế nào khi người dùng chuyển giữa các ký hiệu liên tiếp.
\end{itemize}

\begin{table}[H]
\centering
\caption{Đánh giá định tính mô-đun nhận diện qua webcam}
\label{tab:recognition-qualitative}
\begin{tabularx}{\textwidth}{|X|p{3.2cm}|X|}
\hline
\textbf{Kịch bản thử nghiệm} & \textbf{Kết quả} & \textbf{Nhận xét} \\
\hline
Ký hiệu từ phổ biến như \texttt{hello}, \texttt{drink} & Hoạt động được trong demo & Các nhãn phổ biến có thể xuất hiện ổn định hơn khi tay nằm rõ trong khung hình và ánh sáng đủ. \\
\hline
Ký hiệu liên tiếp nhiều từ & Cần xác nhận từng token & Phiên bản hiện tại phù hợp hơn với cơ chế chọn token/top-1 vào buffer, chưa phải nhận diện câu liên tục hoàn toàn tự động. \\
\hline
Điều kiện ánh sáng yếu & Dễ giảm ổn định & MediaPipe có thể mất landmarks tay hoặc pose, kéo theo dự đoán TFLite kém tin cậy. \\
\hline
Tốc độ ký hiệu nhanh & Có thể bỏ sót động tác & Cửa sổ trượt cần đủ tối thiểu 16 frame; thao tác quá nhanh làm chuỗi landmarks thiếu ngữ cảnh. \\
\hline
Khoảng cách xa camera & Dễ mất chi tiết bàn tay & Khi bàn tay quá nhỏ hoặc bị che, landmarks không ổn định và xác suất dự đoán dao động mạnh. \\
\hline
\end{tabularx}
\end{table}

\subsubsection{Đánh giá định tính và phân tích lỗi}

Bên cạnh các quan sát trực tiếp, phần nhận diện cần có thêm phân tích lỗi theo các hướng:
\begin{itemize}
    \item \textbf{nhầm lớp do động tác gần giống}: nhiều từ khác nhau chỉ lệch hướng cổ tay hoặc biên
    độ chuyển động;
    \item \textbf{nhầm do điều kiện ghi hình}: ánh sáng yếu, nền lộn xộn, tay che nhau hoặc ra ngoài
    khung hình;
    \item \textbf{nhầm do sự khác biệt signer}: mô hình được huấn luyện trên signer trong dữ liệu
    Kaggle, có thể chưa khái quát tốt cho người dùng thực tế;
    \item \textbf{giới hạn 250 từ}: nếu người dùng thực hiện ký hiệu không nằm trong 250 lớp,
    mô hình sẽ dự đoán sai.
\end{itemize}

\subsubsection{Lưu ý khi diễn giải kết quả}

Ngay cả khi mô hình nhận diện hoạt động tốt, vẫn cần nhấn mạnh rằng:
\begin{itemize}
    \item đây mới là chất lượng ở tầng phân loại gloss, chưa phải chất lượng câu tiếng Việt hay tiếng
    Lào cuối cùng;
    \item mô hình được huấn luyện trên dữ liệu ASL, nên kết quả tốt không đồng nghĩa mô hình
    đã ``hiểu'' trọn vẹn VSL thật trong thực địa;
    \item cần tách bạch rõ \textit{recognition accuracy}, \textit{sentence quality}
    và \textit{translation quality}, thay vì gộp chung thành một đánh giá mơ hồ.
\end{itemize}

Trình bày như vậy vừa chính xác về mặt học thuật, vừa giúp người đọc hiểu đúng đóng góp thật của mô
hình nhận diện trong toàn bộ hệ thống.

